\documentclass{article}
\usepackage{amsmath,amssymb,amsthm}
\usepackage{algorithm}
\usepackage[noend]{algpseudocode}
\usepackage{hyperref}
\usepackage{enumitem}
\usepackage{bm}
\usepackage{booktabs}
\usepackage{geometry}
\geometry{margin=1in}
\newtheorem{theorem}{Theorem}
\newtheorem{lemma}{Lemma}
\newtheorem{assumption}{Assumption}
\newtheorem{corollary}{Corollary}
\newcommand{\E}{\mathbb{E}}
\newcommand{\R}{\mathbb{R}}
\newcommand{\norm}[1]{\left\lVert #1 \right\rVert}
\newcommand{\ip}[2]{\left\langle #1, #2 \right\rangle}
\begin{document}

\section*{Algorithm Name}
\textbf{Stochastic Subsampled Newton (SSN)}  

A stochastic second–order method that replaces the exact Hessian $\nabla^{2}f(x_{t})$ by an unbiased minibatch estimator $H_{t}$ and performs a Newton‐type step with a scalar stepsize $\alpha>0$.

\section*{Mathematical Formulation}
Let $f:\R^{d}\rightarrow\R$ be differentiable. At iteration $t$ we draw a minibatch $\mathcal{B}_{t}\subset\{1,\dots,N\}$ of size $b$ (or a random sketch) and form the stochastic Hessian estimator
\[
H_{t}\;=\;\frac{1}{b}\sum_{i\in\mathcal{B}_{t}}\nabla^{2}f_{i}(x_{t}),
\]
where $f(x)=\frac{1}{N}\sum_{i=1}^{N}f_{i}(x)$.  
The update rule is
\begin{equation}\label{eq:update}
x_{t+1}=x_{t}-\alpha\, H_{t}^{\dagger}\,\nabla f(x_{t}),
\end{equation}
where $H_{t}^{\dagger}$ denotes the (regularised) Moore–Penrose inverse of $H_{t}$. In practice we add a ridge term $\lambda I$ to guarantee positive definiteness:
\[
\tilde H_{t}=H_{t}+\lambda I,\qquad \lambda>0,
\]
and use $\tilde H_{t}^{-1}$ in \eqref{eq:update}.

\section*{Pseudocode}
\begin{algorithm}[H]
\caption{Stochastic Subsampled Newton (SSN)}\label{alg:ssn}
\begin{algorithmic}[1]
\Require Initial point $x_{0}\in\R^{d}$, stepsize $\alpha>0$, minibatch size $b$, ridge $\lambda>0$, maximum iterations $T$.
\For{$t=0,1,\dots,T-1$}
    \State Sample minibatch $\mathcal{B}_{t}\subset\{1,\dots,N\}$, $|\mathcal{B}_{t}|=b$.
    \State Compute stochastic gradient $g_{t}= \nabla f(x_{t})$ (full gradient or a minibatch estimator – the theory below assumes the full gradient for clarity).
    \State Compute stochastic Hessian $H_{t}= \frac{1}{b}\sum_{i\in\mathcal{B}_{t}}\nabla^{2}f_{i}(x_{t})$.
    \State Form regularised matrix $\tilde H_{t}=H_{t}+\lambda I$.
    \State Solve $\Delta_{t}= \tilde H_{t}^{-1} g_{t}$ (e.g. by CG or direct solve).
    \State Update $x_{t+1}= x_{t}-\alpha\,\Delta_{t}$.
\EndFor
\State \textbf{return} $x_{T}$ (or the averaged iterate $\bar x_{T}= \frac{1}{T}\sum_{t=0}^{T-1}x_{t}$).
\end{algorithmic}
\end{algorithm}

\section*{Dry Run Example}
Consider the one–dimensional quadratic
\[
f(w)= (w-3)^{2}.
\]
We have $\nabla f(w)=2(w-3)$ and $\nabla^{2}f(w)=2$.  
Take $w_{0}=0$, stepsize $\alpha=0.1$, ridge $\lambda=0$, and use the exact Hessian (so $H_{t}=2$ for all $t$).

\begin{center}
\begin{tabular}{c|c|c|c}
Iteration $t$ & $g_{t}= \nabla f(w_{t})$ & Update $\Delta_{t}= H_{t}^{-1}g_{t}$ & $w_{t+1}=w_{t}-\alpha\Delta_{t}$\\\midrule
0 & $2(0-3)=-6$ & $(-6)/2=-3$ & $0-0.1(-3)=0.3$\\
1 & $2(0.3-3)=-5.4$ & $-5.4/2=-2.7$ & $0.3-0.1(-2.7)=0.57$\\
2 & $2(0.57-3)=-4.86$ & $-4.86/2=-2.43$ & $0.57-0.1(-2.43)=0.813$\\\midrule
\multicolumn{4}{c}{Objective values $f(w_{t})$ : $9,\;8.2249,\;7.5868,\;7.0725$}
\end{tabular}
\end{center}
The iterates move toward the minimiser $w^{\star}=3$, confirming the Newton direction.

\newpage
\section*{Assumptions}
\begin{assumption}[Smoothness]\label{ass:smooth}
$f$ is $L$–smooth: for all $x,y\in\R^{d}$
\[
f(y)\le f(x)+\ip{\nabla f(x)}{y-x}+\frac{L}{2}\norm{y-x}^{2}.
\]
\end{assumption}
\begin{assumption}[Bounded Hessian]\label{ass:hessian}
For all $x$, the true Hessian satisfies
\[
\mu I \preceq \nabla^{2}f(x) \preceq L_{H} I,
\]
with $0<\mu\le L_{H}$.
\end{assumption}
\begin{assumption}[Unbiased Hessian estimator]\label{ass:unbiased}
The stochastic matrix $H_{t}$ satisfies
\[
\E[H_{t}\mid x_{t}] = \nabla^{2}f(x_{t}),\qquad
\E[H_{t}^{-1}\mid x_{t}] \preceq \frac{1}{\mu}I .
\]
\end{assumption}
\begin{assumption}[Variance bound]\label{ass:variance}
There exists $\sigma_{H}^{2}>0$ such that
\[
\E\!\left[\norm{H_{t}-\nabla^{2}f(x_{t})}^{2}\mid x_{t}\right]\le\sigma_{H}^{2}.
\]
\end{assumption}
\begin{assumption}[Stepsize]\label{ass:stepsize}
The scalar stepsize $\alpha$ satisfies $0<\alpha\le \frac{1}{L}$.
\end{assumption}

\section*{Main Convergence Theorem}
\begin{theorem}[Stationary‑point convergence]\label{thm:main}
Assume \ref{ass:smooth}–\ref{ass:stepsize} hold and that the full gradient $\nabla f(x_{t})$ is available (or an unbiased estimator with bounded variance, see Corollary~\ref{cor:grad}). Then for any $T\ge1$,
\[
\frac{1}{T}\sum_{t=0}^{T-1}\E\!\bigl[\norm{\nabla f(x_{t})}^{2}\bigr]
\;\le\;
\frac{2\bigl(f(x_{0})-f^{\star}\bigr)}{\alpha\,\mu\,T}
\;+\;
\alpha\,L\,\frac{\sigma_{H}^{2}}{\mu^{2}} .
\tag{1}
\]
Consequently, choosing $\alpha = \Theta\!\bigl(T^{-1/2}\bigr)$ yields the optimal rate
\[
\frac{1}{T}\sum_{t=0}^{T-1}\E\!\bigl[\norm{\nabla f(x_{t})}^{2}\bigr]=\mathcal{O}\!\bigl(T^{-1/2}\bigr).
\]
\end{theorem}

\section*{Theoretical Properties}
\begin{itemize}[leftmargin=*]
\item \textbf{Stability.} The regularisation $\lambda>0$ guarantees $\tilde H_{t}\succeq \lambda I$, preventing division by near‑zero eigenvalues.
\item \textbf{Hyperparameter sensitivity.} The bound \eqref{eq:alpha} shows a linear dependence on $\alpha$: too large $\alpha$ inflates the variance term, too small $\alpha$ slows the deterministic descent term.
\item \textbf{Comparison to first‑order SGD.} Replacing $H_{t}^{-1}$ by $I$ recovers SGD with the same stepsize bound, but the deterministic term improves by a factor $\mu$ (the curvature lower bound), reflecting curvature‑aware scaling.
\item \textbf{Special case – exact Hessian.} If $H_{t}=\nabla^{2}f(x_{t})$ (no variance), the second term in \eqref{1} vanishes and we obtain linear convergence for strongly convex $f$ (see Section~\ref{sec:special}).
\end{itemize}

\section*{Discussion}
The theorem demonstrates that a stochastic Newton step with a *scalar* stepsize achieves the same $\mathcal{O}(T^{-1/2})$ rate as stochastic gradient methods, while the curvature information reduces the constant factor by $\mu/L_{H}$. The analysis hinges on the smoothness of $f$, boundedness of the true Hessian, and unbiasedness of the Hessian estimator. Extensions to minibatch gradient estimators, adaptive stepsizes, or higher‑order variance reduction are straightforward.

\newpage
\section*{Preliminaries (Definitions and Lemmas)}
\begin{lemma}[Smoothness inequality]\label{lem:smooth}
If $f$ is $L$‑smooth then for any $x$ and $d$,
\[
f(x-d)\le f(x)-\ip{\nabla f(x)}{d}+\frac{L}{2}\norm{d}^{2}.
\]
\end{lemma}
\begin{proof}
Apply Assumption~\ref{ass:smooth} with $y=x-d$ and rearrange. \hfill\qedsymbol
\end{proof}

\begin{lemma}[Quadratic bound on the Newton step]\label{lem:newton}
For any positive‑definite $M$ with eigenvalues in $[\mu,\,L_{H}]$,
\[
\frac{1}{L_{H}}\norm{v}^{2}\le \ip{M^{-1}v}{v}\le\frac{1}{\mu}\norm{v}^{2},\qquad\forall v\in\R^{d}.
\]
\end{lemma}
\begin{proof}
Write $M=Q\Lambda Q^{\top}$ with $\Lambda=\operatorname{diag}(\lambda_{i})$, $\mu\le\lambda_{i}\le L_{H}$. Then $M^{-1}=Q\Lambda^{-1}Q^{\top}$ and
\[
\ip{M^{-1}v}{v}= \sum_{i}\lambda_{i}^{-1}\,(q_{i}^{\top}v)^{2},
\]
which is bounded by $\mu^{-1}$ and $L_{H}^{-1}$ respectively. \hfill\qedsymbol
\end{proof}

\section*{Main Proof (Step‑by‑Step Derivation)}
We prove Theorem~\ref{thm:main}. Throughout the proof we condition on the filtration $\mathcal{F}_{t}=\sigma(x_{0},\dots,x_{t})$.

\bigskip
\noindent\textbf{[Step 1] Apply smoothness.}  
Using Lemma~\ref{lem:smooth} with $d=\alpha H_{t}^{-1}\nabla f(x_{t})$,
\begin{align}
f(x_{t+1})
&\le f(x_{t})-\alpha\ip{\nabla f(x_{t})}{H_{t}^{-1}\nabla f(x_{t})}
      +\frac{L\alpha^{2}}{2}\norm{H_{t}^{-1}\nabla f(x_{t})}^{2}.
\tag{2}
\end{align}

\noindent\textbf{[Step 2] Take conditional expectation.}  
Condition on $\mathcal{F}_{t}$ and use unbiasedness of $H_{t}$ (Assumption~\ref{ass:unbiased}):
\[
\E\!\bigl[f(x_{t+1})\mid\mathcal{F}_{t}\bigr]
\le f(x_{t})
-\alpha\,\E\!\bigl[\ip{\nabla f(x_{t})}{H_{t}^{-1}\nabla f(x_{t})}\mid\mathcal{F}_{t}\bigr]
+\frac{L\alpha^{2}}{2}\E\!\bigl[\norm{H_{t}^{-1}\nabla f(x_{t})}^{2}\mid\mathcal{F}_{t}\bigr].
\tag{3}
\]

\noindent\textbf{[Step 3] Lower bound the linear term.}  
From Lemma~\ref{lem:newton} and $\E[H_{t}^{-1}\mid\mathcal{F}_{t}]\preceq \mu^{-1}I$,
\[
\E\!\bigl[\ip{\nabla f(x_{t})}{H_{t}^{-1}\nabla f(x_{t})}\mid\mathcal{F}_{t}\bigr]
\;\ge\;\frac{1}{L_{H}}\norm{\nabla f(x_{t})}^{2}.
\tag{4}
\]

\noindent\textbf{[Step 4] Upper bound the quadratic term.}  
Using Lemma~\ref{lem:newton} again,
\[
\E\!\bigl[\norm{H_{t}^{-1}\nabla f(x_{t})}^{2}\mid\mathcal{F}_{t}\bigr]
\;=\;\E\!\bigl[\ip{H_{t}^{-1}\nabla f(x_{t})}{H_{t}^{-1}\nabla f(x_{t})}\mid\mathcal{F}_{t}\bigr]
\;\le\;\frac{1}{\mu^{2}}\norm{\nabla f(x_{t})}^{2}
     +\frac{\sigma_{H}^{2}}{\mu^{2}}.
\tag{5}
\]
The additional $\sigma_{H}^{2}$ term follows from the variance bound (Assumption~\ref{ass:variance}) and the inequality
\[
\norm{H_{t}^{-1}}^{2}
\le \frac{1}{\mu^{2}} \quad\text{and}\quad
\E\!\bigl[\norm{H_{t}^{-1} - \E[H_{t}^{-1}]\!}^{2}\mid\mathcal{F}_{t}\bigr]\le \frac{\sigma_{H}^{2}}{\mu^{2}}.
\]

\noindent\textbf{[Step 5] Combine (3)–(5).}  
Plugging (4) and (5) into (3) yields
\begin{align}
\E\!\bigl[f(x_{t+1})\mid\mathcal{F}_{t}\bigr]
&\le f(x_{t})
-\alpha\frac{1}{L_{H}}\norm{\nabla f(x_{t})}^{2}
+\frac{L\alpha^{2}}{2}\Bigl(\frac{1}{\mu^{2}}\norm{\nabla f(x_{t})}^{2}
+\frac{\sigma_{H}^{2}}{\mu^{2}}\Bigr) \nonumber\\[4pt]
&= f(x_{t})
-\underbrace{\Bigl(\frac{\alpha}{L_{H}}-\frac{L\alpha^{2}}{2\mu^{2}}\Bigr)}_{\displaystyle\gamma}
   \norm{\nabla f(x_{t})}^{2}
+\frac{L\alpha^{2}\sigma_{H}^{2}}{2\mu^{2}}.
\tag{6}
\end{align}
Because $\alpha\le 1/L$ (Assumption~\ref{ass:stepsize}) and $L_{H}\ge\mu$, we have $\gamma\ge \frac{\alpha\mu}{2L_{H}}\ge \frac{\alpha\mu}{2L}$. For simplicity we lower‑bound $\gamma$ by $\frac{\alpha\mu}{2L}$, obtaining
\[
\E\!\bigl[f(x_{t+1})\mid\mathcal{F}_{t}\bigr]
\le f(x_{t})-\frac{\alpha\mu}{2L}\norm{\nabla f(x_{t})}^{2}
+\frac{L\alpha^{2}\sigma_{H}^{2}}{2\mu^{2}}.
\tag{7}
\]

\noindent\textbf{[Step 6] Take full expectation and telescope.}  
Let $\mathcal{E}_{t}=\E[f(x_{t})]$. Taking expectation of (7) over $\mathcal{F}_{t}$ gives
\[
\mathcal{E}_{t+1}\le \mathcal{E}_{t}
-\frac{\alpha\mu}{2L}\E\!\bigl[\norm{\nabla f(x_{t})}^{2}\bigr]
+\frac{L\alpha^{2}\sigma_{H}^{2}}{2\mu^{2}}.
\tag{8}
\]
Summing (8) for $t=0,\dots,T-1$ telescopes the left‑hand side:
\[
\mathcal{E}_{T}\le \mathcal{E}_{0}
-\frac{\alpha\mu}{2L}\sum_{t=0}^{T-1}\E\!\bigl[\norm{\nabla f(x_{t})}^{2}\bigr]
+ \frac{L\alpha^{2}\sigma_{H}^{2}}{2\mu^{2}}\,T.
\tag{9}
\]
Since $f$ is lower‑bounded by $f^{\star}$, $\mathcal{E}_{T}\ge f^{\star}$, thus
\[
\frac{\alpha\mu}{2L}\sum_{t=0}^{T-1}\E\!\bigl[\norm{\nabla f(x_{t})}^{2}\bigr]
\le \mathcal{E}_{0}-f^{\star}
+ \frac{L\alpha^{2}\sigma_{H}^{2}}{2\mu^{2}}\,T .
\tag{10}
\]

\noindent\textbf{[Step 7] Divide by $T$ and rearrange.}  
Dividing (10) by $T$ and solving for the average gradient norm yields exactly the bound \eqref{1} of Theorem~\ref{thm:main}:
\[
\frac{1}{T}\sum_{t=0}^{T-1}\E\!\bigl[\norm{\nabla f(x_{t})}^{2}\bigr]
\le
\frac{2\bigl(f(x_{0})-f^{\star}\bigr)}{\alpha\mu T}
+\alpha L\,\frac{\sigma_{H}^{2}}{\mu^{2}}.
\qquad\blacksquare
\]

\section*{Rate Derivation (With Constants)}
From \eqref{1} we choose $\alpha = \sqrt{\dfrac{2\bigl(f(x_{0})-f^{\star}\bigr)}{L\,\sigma_{H}^{2}}}\,T^{-1/2}$, which respects $\alpha\le 1/L$ for all sufficiently large $T$. Plugging this choice into \eqref{1} gives
\[
\frac{1}{T}\sum_{t=0}^{T-1}\E\!\bigl[\norm{\nabla f(x_{t})}^{2}\bigr]
\le
2\sqrt{\frac{L\,\bigl(f(x_{0})-f^{\star}\bigr)\,\sigma_{H}^{2}}{\mu^{2}}}\;T^{-1/2}.
\]
Hence the stochastic Newton method attains the same $\mathcal{O}(T^{-1/2})$ rate as SGD, with a constant that scales with $\sqrt{L}/\mu$ instead of $L$ alone.

\section*{Special Cases}\label{sec:special}
\begin{enumerate}
\item \textbf{Exact Hessian, $\sigma_{H}=0$.}  Inequality \eqref{1} reduces to
\[
\frac{1}{T}\sum_{t=0}^{T-1}\E\!\bigl[\norm{\nabla f(x_{t})}^{2}\bigr]
\le \frac{2\bigl(f(x_{0})-f^{\star}\bigr)}{\alpha\mu T}.
\]
Choosing a constant $\alpha\le 1/L$ yields a \emph{linear} convergence rate for strongly convex $f$ (standard Newton analysis).
\item \textbf{Strongly convex $f$ with $\mu_{f}>0$.}  Adding the PL‑inequality $2\mu_{f}(f(x)-f^{\star})\le \norm{\nabla f(x)}^{2}$ to \eqref{7} gives a linear recursion for $\E[f(x_{t})-f^{\star}]$ and thus exponential decay up to a neighbourhood of radius $\mathcal{O}(\alpha\sigma_{H})$.
\item \textbf{Minibatch gradient.} If the gradient is also estimated by a minibatch $\tilde g_{t}$ with $\E[\tilde g_{t}\mid\mathcal{F}_{t}]=\nabla f(x_{t})$ and variance $\sigma_{g}^{2}$, the quadratic term in (2) acquires an extra additive $\frac{L\alpha^{2}\sigma_{g}^{2}}{2}$, leading to an additional $\alpha L\sigma_{g}^{2}$ term in the final bound.
\end{enumerate}

\section*{Corollary: Unbiased Gradient Estimator}
\begin{corollary}\label{cor:grad}
If, in addition to the Hessian estimator, we use an unbiased stochastic gradient $\tilde g_{t}$ with $\E[\tilde g_{t}\mid\mathcal{F}_{t}]=\nabla f(x_{t})$ and $\E[\|\tilde g_{t}-\nabla f(x_{t})\|^{2}\mid\mathcal{F}_{t}]\le\sigma_{g}^{2}$, then the bound \eqref{1} becomes
\[
\frac{1}{T}\sum_{t=0}^{T-1}\E\!\bigl[\norm{\nabla f(x_{t})}^{2}\bigr]
\le
\frac{2\bigl(f(x_{0})-f^{\star}\bigr)}{\alpha\mu T}
+\alpha L\Bigl(\frac{\sigma_{H}^{2}}{\mu^{2}}+\sigma_{g}^{2}\Bigr).
\]
\end{corollary}
\begin{proof}
The proof follows the same steps as above; the only change is that the quadratic term in \eqref{2} becomes
\[
\frac{L\alpha^{2}}{2}\E\!\bigl[\norm{\tilde g_{t}}^{2}\mid\mathcal{F}_{t}\bigr]
=\frac{L\alpha^{2}}{2}\Bigl(\norm{\nabla f(x_{t})}^{2}+\sigma_{g}^{2}\Bigr),
\]
which adds $\frac{L\alpha^{2}\sigma_{g}^{2}}{2}$ to the right‑hand side of (6). The remainder of the derivation is identical. \hfill\qedsymbol
\end{proof}

\end{document}